% !TeX root = ../main.tex
% Das Quellenverzeichnis ist ein Unterabschnitt des Kapitels "Verzeichnisse".
%
% Die thebibliography-Umgebung von scrreprt setzt ihre Überschrift über
% \chapter*, was einen Seitenumbruch auslöst — die Liste würde von ihrer
% Überschrift getrennt. Innerhalb der Gruppe wird \chapter* daher zum no-op;
% die Überschrift setzt der \section-Befehl, wodurch Nummerierung,
% Inhaltsverzeichnis-Eintrag und PDF-Bookmark regulär entstehen.
\section{Quellenverzeichnis}
\label{sec:quellenverzeichnis}

\begingroup
\makeatletter
\renewcommand{\chapter}{\@ifstar\@gobble\@gobble}
\makeatother
% ieeetr-hm.bst: Ableitung von ieeetr.bst, setzt Titel in Anführungszeichen
% UND kursiv und behält die Groß-/Kleinschreibung aus literatur.bib bei.
% Nummerierung nach Reihenfolge der ersten Nennung im Text.
\bibliographystyle{ieeetr-hm}
\bibliography{literatur}
\endgroup
